\chapter{Background and related work}
\label{sec:related}

\section{Sign language processing tasks and how they are evaluated}
Sign language processing covers four related tasks that differ in their input, output and metric
(Table~\ref{tab:tasks}); a sign-generation system touches all four, because its own output can be checked only with the
recognition and translation tasks.

\begin{itemize}[leftmargin=*,itemsep=2pt]
  \item \textbf{Isolated sign recognition} is a classification task: one clip, one sign, scored by top-1 and top-5
        accuracy. WLASL \cite{wlasl}, MS-ASL \cite{msasl}, ASL Citizen \cite{aslcitizen}, KArSL \cite{karsl} and AUTSL
        \cite{autsl} are benchmarks of this kind, and their recordings are the raw material of sign lexicons.
  \item \textbf{Continuous sign recognition} maps a signed sentence to its \emph{gloss} sequence, the written labels of
        the signs in signing order. Models are trained with connectionist temporal classification (CTC) \cite{ctc},
        which needs only the gloss order, not the boundaries, and are scored by gloss word error rate (WER)
        \cite{koller2015,kollersurvey}. The best systems now reach low error rates on controlled data: the winning
        skeleton-based model of the ICCV 2025 multimodal sign language recognition challenge \cite{mslr2025,signeval2025}
        was trained on the Saudi Isharah dataset \cite{isharah}.
  \item \textbf{Sign language translation} maps a signed sentence to spoken-language text, with or without glosses as an
        intermediate step \cite{camgoz2018,signlanguagetransformers}. Because the target is free text, it is scored with
        machine translation metrics: BLEU \cite{bleu}, ROUGE \cite{rouge}, chrF \cite{chrf}, and learned metrics such as
        BLEURT \cite{bleurt} or BERTScore \cite{bertscore}. Large weakly aligned corpora such as YouTube-ASL
        \cite{youtubeasl} and How2Sign \cite{how2sign} pair sentences with video but have no glosses; pose-based
        translation on them has been studied in ablation \cite{t5slt}.
  \item \textbf{Sign language production} (SLP) is the reverse: text or gloss to a pose sequence or video. Progressive
        Transformers \cite{progressivetransformers} generate continuous skeleton sequences and introduced
        \emph{back-translation} evaluation: a translation model reads the generated pose, and its output is scored with
        BLEU and ROUGE against the source. Pose-level metrics such as DTW-MJE \cite{dtwmje} compare generated and
        reference joints. Text2Sign \cite{text2sign} combines translation with a motion graph and video generation, and
        SignLLM \cite{signllm} trains large multilingual production models. Surveys \cite{slpreview,slpsurvey,slpathways}
        and the SLRTP 2025 production challenge \cite{slrtp2025} report that generated signing still loses hand detail
        and that automatic metrics only partly agree with Deaf judges.
\end{itemize}

\begin{table}[ht]
\centering\small
\caption{Sign language processing tasks, their standard metrics, and their role in Isharati.}
\label{tab:tasks}
\begin{tabularx}{\textwidth}{L{3.0cm} L{2.6cm} L{3.3cm} Y}
\toprule
Task & Input $\rightarrow$ output & Standard metrics & Role in Isharati \\
\midrule
Isolated recognition & clip $\rightarrow$ one sign & top-1 / top-5 accuracy & source of lexicon signs; sign-level back-translation \\
Continuous recognition & sentence video $\rightarrow$ glosses & gloss WER & sentence-level back-translation \\
Translation & sentence video $\rightarrow$ text & BLEU, ROUGE, chrF, BLEURT & meaning-level back-translation \\
Text-to-gloss & text $\rightarrow$ glosses & gloss WER, BLEU & glossing step, evaluated against ArabSign \\
Production & text / gloss $\rightarrow$ pose & back-translation BLEU / WER, DTW-MJE, human rating & the output of the system \\
\bottomrule
\end{tabularx}
\end{table}

Concatenating recorded dictionary signs, the production approach taken here, gives real hand shapes for every sign at
the cost of transitions between signs; it suits a domain where the vocabulary must be exact. Glosses matter for it
twice: the answer is translated into glosses, and each gloss is the key into the lexicon.

\section{Religious texts in sign language}
Religious texts are among the best-resourced content in sign language, but almost entirely for one tradition. Complete
Bibles exist in American Sign Language \cite{aslv,jwaslbible}; the Deaf Bible Society distributes Scripture in 92 sign
languages \cite{deafbible}; and the JWSign corpus already makes 2,530 hours of Bible translations in 98 sign languages,
by more than 1,500 signers, available for machine learning \cite{jwsign}, followed by openly licensed sign Bibles for
model training \cite{osb}. For Islam, the machine-learning resources we found are small and narrow: image recognition of
the isolated letters that open some surahs \cite{quransign1,quransign2}, and an Arabic Sign Language production dataset
of Friday-sermon content with 262 videos \cite{quransign3}, used to train a speech-to-pose model \cite{quransign4},
which is the closest prior work to ours. We found no aligned Qur'an or hadith corpus in any sign language, no lexicon of
Islamic signs across sign languages, and no system that answers religious questions in sign language from authenticated
sources.

\section{Arabic Sign Language}
Arabic Sign Language is a family of national sign languages; the League of Arab States promoted a unified dictionary
\cite{arslunified}. Datasets include KArSL (502 isolated signs) \cite{karsl}, ArabSign (continuous sentences with gloss
annotation) \cite{arabsign} and the Saudi Isharah dataset \cite{isharah}. Recent Arabic systems translate text to gloss
\cite{arglosstrans}, translate Saudi Sign Language with T5 \cite{saudit5}, build parallel corpora for Syrian ArSL
\cite{syrisign}, or drive 3D avatars \cite{algerianavatar,arabicavatar2024,arabic3davatar2022}. None of these targets
religious questions with a grounding guarantee, and none covers several sign languages from the same knowledge base.

\section{Other sign languages used here}
For ASL: WLASL \cite{wlasl}, MS-ASL \cite{msasl}, ASL Citizen \cite{aslcitizen}, How2Sign \cite{how2sign}, the Boston
ASLLVD \cite{boston} and YouTube-ASL \cite{youtubeasl,youtubeaslkp}. For Turkish: the official T\.{I}D dictionary
\cite{tidsozluk} and AUTSL \cite{autsl}. For Indian Sign Language, used for Urdu because large public Pakistan Sign
Language datasets are lacking \cite{psl}: CISLR \cite{cislr}, iSign \cite{isign} and the WSLP 2025 shared task
\cite{wslp,wslpislword}. Signs are found in weakly labelled sentence video by sign spotting, using subtitles or mouthings
as weak labels \cite{bsl1k,momeni2022}. Keypoints throughout come from MediaPipe \cite{mediapipe} Holistic
\cite{mediapipeholistic,blazepose}.

\section{Grounded generation}
Retrieval-augmented generation \cite{rag} conditions a language model on retrieved passages. Isharati combines BM25
\cite{bm25} with multilingual E5 embeddings \cite{e5,me5} through reciprocal rank fusion \cite{rrf}, and adds checks after
generation, because retrieval alone does not stop a model from adding facts.
